% sections/01-introduction.tex

\section{Introduction}\label{sec:introduction}

The Four Colour Theorem (4CT), that every planar graph admits a proper
$4$-coloring, was first proved by Appel and Haken~\cite{appel1977}
using an unavoidable set of 1,476 reducible configurations, each
verified by computer. A cleaner proof by Robertson, Sanders, Seymour,
and Thomas~\cite{robertson1997} reduced the set to 633 configurations,
but the essential reliance on exhaustive computation remains. In the
half-century since, no ``human-readable'' proof has emerged.

A fundamentally different strategy proceeds \emph{constructively}: begin
with the easy Five Colour Theorem, then systematically eliminate the
fifth color via Kempe chain swaps. Kempe's 1879 ``proof'' used exactly
this strategy; Heawood's 1890 counterexample showed only that a
particular swap could merge two chains, not that the strategy itself
fails. Mathewson~\cite{mathewson2025constructive} revived and formalized this line of
attack. An inductive proof architecture based on vertex removal and
lifting achieves a nearly complete proof: six lemmas are proved,
covering all cases except when the removed vertex $v$ has degree 4 or 5
and all five colors appear among its neighbors. In this critical case,
one must show that the 5-coloring of $G - v$ can be reconfigured,
via a sequence of Kempe swaps, into a coloring where two neighbors of
$v$ share a color, freeing a color for $v$.

\subsection{The Avoidance Problem}

The central obstacle is \emph{avoidance}: during reconfiguration, one
must avoid Kempe swaps that merge the two chains incident to $v$, as
such merges destroy the degree of freedom needed to recolor $v$.
Conjecture~5.5 of the constructive architecture asserted that
BFS-optimal paths in the reconfiguration graph $\R(G - v, 5)$ always
avoid such unsafe swaps. Exhaustive enumeration at $n = 9$ disproves this
conjecture: for the triangulations $T_{9,25}$ and $T_{9,35}$, \emph{all}
optimal-length paths from certain 5-colorings to the nearest 4-coloring
pass through at least one unsafe swap~\cite{mathewson2025}.

The counterexamples refute only the optimality claim, not the
constructive strategy itself: in every case, \emph{safe} non-optimal paths
exist, requiring one or two extra steps. This led to the revised
Conjecture~5.5$'$ (Safe Path Existence): for every 5-coloring of a
planar triangulation and every vertex $v$, there exists a
reconfiguration path to a 4-coloring that avoids all swaps merging
chains incident to $v$.

\subsection{Why Physics?}

What distinguishes safe paths from unsafe ones? If we assign each
coloring an energy via some functional, the reconfiguration graph becomes
a discrete energy surface, and safe versus unsafe paths trace distinct
trajectories. We argue that anti-ferromagnetic spin models and disordered
systems provide the discriminating tools.

A proper $k$-coloring of a
graph $G$ is precisely a ground state of the anti-ferromagnetic Potts
model on $G$ at zero temperature. The chromatic polynomial $P(G, k)$ is
the zero-temperature partition function. Kempe chain reconfiguration is
the natural dynamics of this system: a Kempe $(a,b)$-swap flips a
connected cluster of spins in the $(a,b)$-subgraph, which is exactly
the Swendsen-Wang cluster update restricted to two colors.

We pursue a geometric analogy by embedding the four colors at the
vertices of a regular tetrahedron in $\mathbb{R}^3$, with the fifth
color at the centroid. In this embedding, each Kempe swap becomes a
$180^\circ$ rotation in $\mathrm{SO}(3)$, and the symmetry group of the
tetrahedron maps isomorphically to $S_4$, so that reconfiguration
can be viewed metaphorically as a walk through a space of
$\mathrm{SO}(3)$-valued fields on the graph. This picture is
reminiscent of Kauffman's reformulation via the Penrose
evaluation~\cite{kauffman1990} and of topological quantum field theory
(TQFT), though making these connections precise remains an open
problem.

\subsection{Contributions}

In this paper we:
\begin{enumerate}[label=(\roman*)]
  \item Define eight energy functionals for the reconfiguration graph,
    drawn from anti-ferromagnetic Potts models, polyhedral covariance
    (Magic Gems~\cite{mathewson2025}), electrostatic charge flow,
    protein folding landscapes, and surface tension
    (Section~\ref{sec:methods}).
  \item Evaluate these functionals on the first known counterexamples to
    Conjecture~5.5 at $n = 9$ (Section~\ref{sec:results}).
  \item Identify, as computational signatures observed at $n = 9$, three
    candidate discriminators: local entropy, Magic Gem energy
    trajectories, and surface tension rigidity.
  \item Formalize the \emph{Surface Tension Rigidity Conjecture}: chains
    with zero-variance normalized surface tension are merge-prone.
    We note that this conjecture was subsequently falsified at
    $n \geq 10$.
  \item Discuss connections to the Penrose evaluation, sheaf cohomology,
    chromatic polynomials, and the discharging method
    (Section~\ref{sec:discussion}).
\end{enumerate}

\subsection{Outline}

Section~\ref{sec:background} reviews Kempe chains, the constructive 4CT
architecture, the reconfiguration graph, and the relevant physical
models. Section~\ref{sec:methods} defines all eight energy functionals
with the tetrahedral color embedding. Section~\ref{sec:results} presents
computational results on the $n = 9$ counterexamples.
Section~\ref{sec:discussion} explores speculative connections to TQFT,
sheaf cohomology, chromatic polynomials, and discharging.
Section~\ref{sec:conclusion} summarizes and poses open questions.
